\section{Self-assessment}
\label{sec:self}

This section is written to be useful rather than flattering. It records what was actually finished,
what was deliberately left out, where the shipped code disagrees with the plan it was built from, and
what a final review pass found wrong late enough to be worth reporting.

\subsection{What was completed}

Every capability described in the original concept (\code{docs/idea.md}) is implemented and reachable
from the UI. The knowledge base --- branches, folders, subfolders and Markdown posts, with create,
rename, move, reorder and cascading delete --- works, as do wiki-links, backlinks, related posts,
tags, per-post status, favourites, external references, the two-scope dictionary, and FTS4 search
across titles, content, tags, dictionary terms and resources. The AI assistant runs all four modes,
reads the library only through read-only context tools, and can propose but never perform a change.
Backup, import, settings and theming are done. Fourteen destinations --- five tabs and nine pushed
screens --- are wired into a single navigation graph, and every list screen binds the same loading /
empty / error state helper.

Three features were listed in the plan as explicitly \emph{out} of scope for version~1 and were built
anyway, once the graded core was safe:

\begin{itemize}
  \item the force-directed \textbf{knowledge graph}, with a whole-library view and a one-hop
        neighbourhood view, laid out by the project's own solver rather than a library;
  \item \textbf{streamed AI replies} over Server-Sent Events, including the awkward part --- pulling a
        readable answer out of a JSON envelope that has only partly arrived;
  \item \textbf{in-app YouTube playback} in a locked-down bottom sheet, instead of throwing the user
        out to another app.
\end{itemize}

\subsection{What was deliberately not built}

None of these are accidents; each was a scope decision, and each costs something.

\begin{itemize}
  \item \textbf{Posts are Markdown text only.} No images, no file attachments, no rich-text editing.
        This keeps the backup file a single self-contained JSON document with no binary side-car, and
        it keeps the editor honest, but it does rule out the kind of note that is mostly a diagram.
  \item \textbf{Videos are not embedded in the article flow.} A YouTube URL becomes a resource card
        below the article; tapping it opens an in-app player. The concept asked for embedding
        \emph{inside} the post, which would have meant a custom Markwon plugin.
  \item \textbf{Drag and drop only reorders siblings.} Moving an item into a different container is a
        picker dialog, not a drag. Cross-container drag in a tree that is rendered one level at a time
        needs auto-scroll, hover-to-open and a drop-target model; the dialog does the same job.
  \item \textbf{The review screen edits titles and content, not metadata.} Proposed tags, status,
        references and dictionary entries are now \emph{shown} on the row and can be deselected
        wholesale, but they cannot be edited in place before applying.
  \item \textbf{The assistant can delete posts, never containers.} There is no
        \code{delete\_branch} or \code{delete\_folder} operation, by design --- deletion cascades, and
        no undo can bring a subtree back.
  \item \textbf{Duplicate detection in Organize mode rests on the model.} There is no on-device
        similarity check and no action type for merging two posts.
  \item \textbf{Tablet support is three dimension overrides,} not a layout. Gutters, chat-bubble width
        and card width change above 600\,dp; there is no two-pane list-detail arrangement, and this
        was never verified on real tablet hardware --- only on a resized emulator.
  \item \textbf{R8 is off in the release build.} Gson reflects over the AI protocol classes and the
        network DTOs, and a wrong keep rule there fails \emph{silently}: every reply would deserialize
        to nulls and render as ordinary prose. The keep rules are written and correct, but the flag
        stays off, because a larger APK is not graded and a silently broken assistant is.
\end{itemize}

\subsection{Known limits of features that do ship}

\begin{itemize}
  \item \textbf{Undo after an AI batch is narrow.} It is in-memory, single-level, and scoped to one
        batch: it restores title, content, status, tags, favourite and parent for the posts the batch
        touched, and it removes what the batch created. Deleted posts stay deleted. Leaving the review
        screen ends the offer.
  \item \textbf{The graph hides some labels rather than crowding them.} Whether labels appear at all
        is decided from average node spacing; on top of that, a label that would land on one already
        drawn is dropped, best-connected nodes winning. A dense cluster therefore reads cleanly but
        does not name every dot --- tapping one always shows its title.
  \item \textbf{Search ranking is hand-rolled.} FTS4 has no \code{bm25}, so results are ordered in
        Kotlin --- posts first, then title-prefix matches, then alphabetically --- and snippets are
        built by hand rather than by SQLite.
  \item \textbf{The assistant needs the local gateway.} Without \code{backend/} running and a DeepSeek
        key, the chat tab shows an offline banner. Everything else in the app is unaffected, which is
        the point of the architecture, but it does mean the signature feature cannot be demonstrated
        from the APK alone.
\end{itemize}

\subsection{Where the code diverged from the plan}

The plan in \code{docs/plan.md} was written before the code and was not retro-fitted; the code is the
truth. The differences worth naming:

\begin{itemize}
  \item Context results are returned to the model as a labelled ordinary turn, not an OpenAI
        \code{tool} role --- DeepSeek rejects a tool message with no matching \code{tool\_call\_id}.
        This also keeps the protocol portable, which was the original goal anyway.
  \item The context loop allows four model calls per message, not three, and a request that still
        needs context after the cap returns a retryable error rather than forcing an answer.
  \item The validator resolves titles created earlier in the same batch, not only rows already in the
        database --- otherwise "create a post, tag it, add a term to it" could never validate, and that
        is the normal shape of a generated learning path.
  \item The backup envelope carries the six knowledge tables and not the settings, and chat history is
        never exported.
  \item \code{posts\_fts} is content-backed, so the apply transaction does not rebuild it; Room's own
        triggers keep it in step with \code{nodes}.
\end{itemize}

\subsection{What a final review pass found}

Before submission the whole codebase was audited against the concept document and the grading
criteria, and each finding re-checked against the code before anything was changed. The audit is the
part of this project we would repeat: it found problems that testing the happy path never would.

The most instructive one was found by running the app rather than reading it. Asked to build a
Distributed Systems learning path, the model returned seven operations, one of which used an
operation name that does not exist in the protocol. The validator correctly rejected it --- and
because every row arrived ticked, one invented name disabled \textbf{Apply} for the entire batch. Six
good changes were unreachable because of one bad line. Unsupported operations now arrive unticked
with a plain explanation on the row, so a model that improvises costs the user that one change.
Figure~\ref{fig:aiflow} describes the intended path; this was the case where the intended path had a
sharp edge in it.

The rest, briefly:

\begin{itemize}
  \item An AI \code{update\_post} carrying tags \emph{replaced} the post's whole tag set, silently
        dropping tags the user had added --- and the review screen never displayed tags at all, so it
        happened unseen. Tags are now merged, and proposed tags and status are shown on the row.
  \item Importing the same backup twice on the merge path created duplicate post titles. Titles are
        how wiki-links and AI actions address posts, so duplicates make resolution arbitrary. Imported
        copies that collide are now given a distinct title.
  \item Removing a post's last tag left the tag row behind forever, listed as "0 posts" and reported
        to the assistant as if it existed. Empty tags are now purged.
  \item The tags screen read a DAO directly from its ViewModel --- the one place in the app that broke
        the layering rule --- and counted posts per tag with one query each. It now goes through the
        repository with a single grouped query.
  \item The in-app privacy text understated what is sent: a titles-only outline of the library
        accompanies every AI request, which the wording did not mention. It does now.
  \item A gateway address typed without a scheme was stored, silently ignored, and reported as
        unreachable forever --- exactly what someone testing on a physical device would type.
  \item Two error-state Retry buttons did nothing, and several lists dropped their scroll position on
        rotation.
\end{itemize}

Two smaller protocol holes were closed at the same time: a reference to a branch or folder could be
used where a post was required, and a reference carrying stray whitespace passed validation and then
failed inside the transaction.

\subsection{If there were more time}

In order of what we would actually do first: a two-pane tablet layout, since it is the one rubric
criterion the app answers weakly; Espresso coverage of the review-and-apply path, which is the
riskiest flow in the app and currently has no automated test; editable metadata in the review screen;
and cross-container drag in Browse. Further out, the interesting question is whether the action
protocol could describe a merge of two posts safely enough to let Organize mode deduplicate on its
own, which is the one thing the concept asked for that the model, not the app, is currently trusted to
do.
